\sgaStableTarget{sga6.e10.para.7-16-2}\textbf{7.16.2.}
We shall not review here, in the case of geometric fibers, the extension of all the properties of specialization homomorphisms proved in the preceding sections for the ``arithmetic'' fibers; they all follow trivially from the latter by passage to the limit. Let us mention only the following interesting consequence of 7.9.4 when @@M0@@ is smooth and proper: @@M1@@ is isomorphic to a quotient of a subgroup of @@M2@@; in particular, its rank is bounded above by that of the latter group. This is therefore a \emph{semicontinuity} result for the rank of the group of cycles modulo numerical equivalence.

\sgaStableTarget{sga6.e10.sec.7.17}\subsection*{7.17. Case of an arbitrary base.}
Now let
@@M3@@
be a flat and quasi-compact morphism of noetherian preschemes, and let @@M4@@ be two points of @@M5@@ such that @@M6@@ is a specialization of @@M7@@. We wish to relate the groups @@M8@@, etc., for the two fibers @@M9@@ and @@M10@@, and for corresponding geometric fibers. For this, note (EGA II 7.1.9) that there exists a morphism
@@M11@@
where @@M12@@ is a trait, and where @@M13@@ carries the closed point @@M14@@ to @@M15@@ and the generic point @@M16@@ to @@M17@@; one may even assume that @@M18@@. Suppose that @@M19@@ is noetherian, as is the case when @@M20@@ is of finite type or @@M21@@ is essentially of finite type over @@M22@@ (EGA IV 1.3.8) (which one may assume when @@M23@@ itself is universally Japanese (EGA @@M24@@ 23.1.1)). Then the specialization homomorphism for @@M25@@ defines a homomorphism
\sgaStableTarget{sga6.e10.eq.7-17-1}@@M26@@
